\section{Os ciclos do TDD}

Segundo Uncle Bob (\citeauthor{artigo:cycles-tdd}, \citeyear{artigo:cycles-tdd},
\citetitle{artigo:cycles-tdd}), o TDD acontece em quatro ciclos aninhados,
cada um com seu próprio relógio.

\begin{center}
\small
\begin{tabular}{@{}l l p{8cm}@{}}
\toprule
Ciclo & Relógio & O que fazer \\
\midrule
Nano & segundos & As três leis: só escrever código de produção com um teste
falhando, só escrever teste suficiente para falhar, só escrever código
suficiente para passar \\
Micro & minutos & Vermelho, verde, refatorar \\
Milli & $\sim$10 min & Testes cada vez mais específicos, código cada vez
mais genérico \\
Primário & $\sim$1 hora & Parar e olhar as fronteiras da arquitetura \\
\bottomrule
\end{tabular}
\end{center}

Os dois primeiros ciclos são os que você sente no dedo: as três leis
mandam a cadência, e o ritmo vermelho, verde, refatorar fecha cada volta.
O ciclo milli é o de generalização: a cada teste novo o teste fica mais
específico e o código de produção, mais genérico. O ciclo primário é o
momento de levantar a cabeça e olhar a arquitetura, por exemplo onde
ficam as bordas do sistema.

Uma extensão natural da ideia de ``código mais genérico'' é a
\citetitle{artigo:tpp} (\citeauthor{artigo:tpp}, \citeyear{artigo:tpp}),
indicada abaixo em \enquote{onde aprofundar}.
